\chapter{Bilder, Tabellen, Formeln}%
%
%
\section{Label und Referenzen}%
\label{sec: label}%
Damit man auf bestimmte Teile der Arbeit verlinken kann, muss man \emph{Label} setzen. Dabei kann man auf verschiedene Dinge zeigen (\emph{referenzieren}), z.B. eine Seite, ein Kapitel, eine Tabelle oder ein Bild. Die Befehle aus Tabelle~\ref{tab: referenzieren} bleiben dabei gleich, nur die Positionierung ändert sich.%

Dabei empfiehlt es sich, den \emph{Labels} sinnvolle Präfixe zu geben. Bilder kann man z.B. mit »\emph{fig:}« und Tabellen mit »\emph{tab:}« benennen. Man beachte die \emph{Tilde (\textasciitilde)} beim anschließenden Referenzieren.%

{\tabulinesep=1.2mm%
\begin{table}[!hbt]%
\caption{Label und Referenzen}%
\label{tab: referenzieren}%
\begin{tabu} to \textwidth {X[l]X[l]X[l]}%
\toprule%
\textbf{Bezeichnung} & \textbf{Latex Code} & \textbf{Ausgabe/\emph{Beschreibung}}\tabularnewline%
\midrule%
Label setzen & \textbackslash label\{chap: vorwort\} & \textit{setzt ein Label mit dem Namen \emph{chap: vorwort}. Dies wurde im Quellcode bereits vorgenommen.}\tabularnewline
Referenz aufrufen & Kapitel\textasciitilde\textbackslash ref\{chap: vorwort\} & Kapitel~\ref{chap: vorwort}\tabularnewline
\bottomrule%
\end{tabu}%
\end{table}%
}
%
\section{Bilder}%
Um Bilder zu platzieren, müssen diese in eine Umgebung geschrieben werden. In dieser Umgebung platziert man dann das Bild, die Positionierung/Größe/Verschiebung des Bildes, die Beschriftung und den Link. In Tabelle~\ref{tab: abbildung} sind die Befehle in einem Beispiel gelistet.%

Für die Positionierung des Bildes hat man in den Optionen die Wahl zwischen \emph{center}, \emph{outer} und \emph{inner}. \emph{Inner} positioniert das Bild auf der inneren Seite bei zweiseitigem Druck (also links bei einseitigem Druck) und \emph{outer} positioniert es auf der äußeren Seite (einseitiger Druck: rechts). \emph{Center} platziert es in der Mitte.%

Die Beschriftung der Bilder kann entweder über oder unter dem Bild angeordnet sein. Die Position hängt dabei davon ab, ob \emph{\textbackslash caption\{\}} über oder unter dem Befehl \emph{\textbackslash includegraphics\{\}} steht.%

Die \emph{Bildunterschrift} inklusive des \emph{Labels} sollten immer ganz am Ende der Bildumgebung stehen (sinnvolle Präfixe beim Benennen nicht vergessen -- siehe Kapitel~\ref{sec: label}), eine Begründung dazu gibt es sogar in wissenschaftlichen Arbeiten \cite[][S.101]{schlosser}! In Abbildung~\ref{fig: stura-logo} sieht man das Ergebnis aller Befehle kombiniert aus Tabelle~\ref{tab: abbildung}.%

{\tabulinesep=1.2mm%
\begin{table}[!hbt]%
\caption{Befehle zum Einfügen einer Abbildung}%
\label{tab: abbildung}%
\begin{tabu} to \textwidth {X[l]X[l]X[l]}%
\toprule%
\textbf{Bezeichnung} & \textbf{Latex Code} & \textbf{Ausgabe/\emph{Beschreibung}}\tabularnewline%
\midrule%
Bildumgebung starten & \textbackslash begin\{figure\}[!hbt] & \textit{Beginnt eine Abbildungsumgebung und versucht diese in der Reihenfolge \emph{hier (!h), auf der Seite unten (b) oder auf der nächsten Seite oben (t)} zu platzieren.}\tabularnewline%
Bild einfügen, skalieren und positionieren & \seqsplit{\textbackslash includegraphics[width=0.5\textbackslash textwidth ,center]\{stura-logo-htw\}} & \emph{fügt das Bild mit dem Namen \emph{stura-logo-htw} ein. Dieses wird auf die Hälfte der Textbreite skaliert und innen auf der Seite eingefügt.}\tabularnewline%
Bildunterschrift hinzufügen & \textbackslash caption\{Logo des STURA der HTW Dresden\} & \emph{fügt eine Bildunterschrift mit der Bezeichnung \emph{Logo des STURA der HTW Dresden} hinzu}\tabularnewline%
Verlinkung hinzufügen & \textbackslash label\{fig: stura-logo\} & \emph{fügt eine Verlinkung mit dem Namen \emph{fig: abbildung} dem Bild hinzu, auf die referenziert werden kann}\tabularnewline%
Bildumgebung schließen & \textbackslash end\{figure\}[!hbt] & \emph{Schließt die Abbildungsumgebung}\tabularnewline%
\bottomrule%
\end{tabu}%
\end{table}%
}%

\begin{figure}[!hbt]%
    \includegraphics[width=0.5\textwidth]{stura-logo-htw.pdf}%
    \caption{Logo des STURA der HTW Dresden}%
    \label{fig: stura-logo}%
\end{figure}%
%
%

\section{Tabellen}%
%
%
\section{Formeln}%
%
%
\section{Literatur}%
\label{sec: literatur}%
Die Literatur schreibt man in die Datei \emph{literatur.bib}, welche relativ zu eurem Template im Verzeichnis \emph{/frontbackmatter/} liegt. Hier kommt euch \emph{TexMaker} zu Hilfe, denn unter \emph{Bibliography -> biblatex} findet man Verschiedene Arten von Literatur, die man mit einem Klick hinzufügen kann.%

Dabei sollte man beachten, dass zwischen den einzelnen Quellen unbedingt eine \textsc{Leerzeile} sein sollte, da es sonst eventuell zu Problemen beim Referenzieren kommt.%

Im ersten Feld gibt man seinen \emph{Key} für diese Quelle an. Der \emph{Key} funktioniert später beim zitieren wie ein \emph{Label}. Man kann damit auf den entsprechenden Eintrag verweisen. Alle anderen Felder sind größtenteils selbsterklärend und können sonst in der Dokumentation für Biblatex nachgelesen werden.%

Anschließend kann man mit der Option \emph{Bibliography -> Clean} alle optionalen Felder automatisch entfernen lassen, sofern die Buchstaben \emph{OPT} am Anfang der Zeile nicht entfernt wurden.%
%
%
\section{Zitieren}%
\label{sec: zitieren}
Nachdem man seine Quellen dem Literaturverzeichnis hinzugefügt hat (\ref{sec: literatur})und diese mit \emph{Keys} versehen hat versehen hat, kann man diese anschließend Zitieren.%

Für indirekte Zitate muss man nur die Quelle am Ende des Absatzes/Satzes angeben. Wie das geht, sieht man in \ref{tab: zitieren}. Bemerkungen kann man einfügen, muss man aber nicht. Als \emph{Notiz nachher} empfiehlt es sich, z.B. eine Seitenzahl anzugeben wie hier am Ende des Satzes zu sehen. \cite[][S.22 ff]{einstein}%

Natürlich zitiere ich hier nicht aus dem Artikel von Albert Einstein, es dient in diesem Zusammenhang lediglich zu Vorführungszwecken.%
{\tabulinesep=1.2mm%
\begin{table}[!hbt]%
\caption{zitieren}%
\label{tab: zitieren}%
\begin{tabu} to \textwidth {X[l]X[l]X[l]}%
\toprule%
\textbf{Bezeichnung} & \textbf{Latex Code} & \textbf{Ausgabe/\emph{Beschreibung}}\tabularnewline%
\midrule%
indirektes Zitat & \textbackslash cite[Notiz vorher][Notiz nachher]\{key\} & \cite[Notiz vorher][Notiz nachher]{einstein}\tabularnewline%
direktes Zitat &%
\begin{minipage}[t]{0.32\textwidth}
	%\begin{lstlisting} - funktioniert nicht im TABU environment! vielleicht tabularx?
	%\blockcquote[vorher][nachher]{key}{Hier steht ein wörtliches Zitat.}	
	%\end{lstlisting}%
\end{minipage}%
	& \blockcquote[vorher][nachher]{einstein}{Hier steht ein wörtliches Zitat.}\tabularnewline%
Auslassungszeichen bzw. Pünktchen &\textbackslash textelp\{\} & \textelp{} \tabularnewline%
\bottomrule%
\end{tabu}%
\end{table}%
}%

Manchmal möchte man unbedingt manuell Anführungszeichen setzen, wie in \ref{tab: quotationmarks} zu sehen. Man beachte, dass man immer ein Zitat öffnen und schließen muss. Auf ein \guillemotright \space sollte immer später ein \guillemotleft \space folgen, da es sonst zu Fehlern beim Kompilieren kommt.\\%
\begin{table}[h]%
\begin{tabu} to \textwidth {X[l]X[l]X[l]}%
\toprule%
\textbf{Tastenkombination} & \textbf{Latex Code} & \textbf{Befehl}\tabularnewline%
\midrule%
SHIFT+F1 & \guillemotright & Anführungszeichen doppelt öffnen\\ %
SHIFT+F2 & \guillemotleft & Anführungszeichen doppelt schließen\\%
SHIFT+F3 & \guilsinglright & Anführungszeichen einfach öffnen\\%
SHIFT+F4 & \guilsinglleft & Anführungszeichen einfach schließen\\%
\bottomrule%
\end{tabu}%
\caption{Manuelle Anführungszeichen}%
\label{tab: quotationmarks}%
\end{table}%

Für wörtliche Zitate nutzt man den Befehl aus \ref{tab: woertlicheszitat}. Mit diesem Befehl werden Zitate entweder in Anführungszeichen gesetzt oder automatisch in Blöcke geschrieben, sobald sie 3 Zeilen überschreiten. Möchte man Anführungszeichen trotz Blockschreibweise haben, so fügt man die Anführungszeichen aus \ref{tab: quotationmarks} hinzu.\\%
\begin{table}[h]%
\begin{tabu} to \textwidth {X[c]X[c]}%
\toprule%
\textbackslash blockcquote[vorher][nachher]\{key\}\{Hier steht ein wörtliches Zitat.\} & \blockcquote[vorher][nachher]{einstein}{Hier steht ein wörtliches Zitat.}\\%
\bottomrule%
\end{tabu}%
\caption{Wörtliches Zitat}%
\label{tab: woertlicheszitat}%
\end{table}%

Ebenso benötigt man häufig die Auslassungspunkte aus Tabelle~\ref{tab: zitieren}, um für den Sinn des Zitates nicht relevante Textstellen zu auszulassen und zu markieren.%

Für Zitate in anderen Sprachen schaut man in die Dokumentation von dem Paket \emph{csquotes}.%

Damit die Literaturreferenzen im Dokument erkannt werden, muss man vor dem Kompilieren den selbst erstellten \emph{User Command} für \emph{biblatex} starten und anschließend die PDF mindestens zwei mal Kompilieren, damit das Ergebnis vollständig sichtbar wird.%
%
%
\section{Fußnoten}%
Fußnoten setzt man wie in \ref{tab: fussnoten} zu sehen.\\%
{\tabulinesep=1.2mm%
\begin{table}[!hbt]%
\caption{Fußnoten}%
\label{tab: fussnoten}%
\begin{tabu} to \textwidth {X[l]X[l]X[l]}%
\toprule%
\textbf{Bezeichnung} & \textbf{Latex Code} & \textbf{Ausgabe/\emph{Beschreibung}}\tabularnewline%
\midrule%
Fußnote setzen Text & \textbackslash footnote\{Fußnotentext\} & Text \footnote{Fußnotentext} \\%
\bottomrule%
\end{tabu}%
\end{table}%
}%
